% ═══════════════════════════════════════════════════════════════
% AB-02a: Ohmsches Gesetz (didaktisch überarbeitet)
% Physik Klasse 8 | Doppelstunde 02a
% R-Berechnung NACH Einführung des Ohmschen Gesetzes!
% ═══════════════════════════════════════════════════════════════

\documentclass[11pt, a4paper]{article}

\usepackage[utf8]{inputenc}
\usepackage[T1]{fontenc}
\usepackage[ngerman]{babel}

\usepackage{helvet}
\renewcommand{\familydefault}{\sfdefault}

\usepackage{microtype}
\usepackage[margin=1.5cm, top=1.8cm, bottom=1.5cm]{geometry}
\usepackage{enumitem}
\usepackage{tabularx}
\usepackage{booktabs}
\usepackage{array}
\usepackage{multirow}
\usepackage{tikz}
\usetikzlibrary{calc}

\usepackage{amsmath, amssymb}
\usepackage{siunitx}
\sisetup{locale = DE, per-mode = symbol}

\usepackage{fancyhdr}
\usepackage{lastpage}
\usepackage{needspace}

\usepackage{xcolor}
\usepackage{tcolorbox}
\tcbuselibrary{skins, breakable}

% Farben
\definecolor{physik}{HTML}{2c5aa0}
\definecolor{hellgrau}{HTML}{F5F5F5}
\definecolor{mittelgrau}{HTML}{888888}
\definecolor{dunkelgrau}{HTML}{444444}
\definecolor{rahmen}{HTML}{CCCCCC}

% Variablen
\newcommand{\fach}{Physik}
\newcommand{\thema}{Das Ohmsche Gesetz}
\newcommand{\klasse}{8}
\newcommand{\maxpunkte}{47}

\colorlet{fachfarbe}{physik}

\newif\ifzeigesternchen
\zeigesternchentrue

% Kopf- und Fußzeile
\pagestyle{fancy}
\fancyhf{}
\renewcommand{\headrulewidth}{0.4pt}
\lhead{\small \fach{} Klasse \klasse{} | \thema}
\rhead{\small Name: \underline{\hspace{3.5cm}}}
\cfoot{\small\textcolor{mittelgrau}{Seite \thepage{} von \pageref{LastPage}}}

% Boxen
\newtcolorbox{merkkasten}[1][]{
    colback=hellgrau,
    colframe=hellgrau,
    boxrule=0pt,
    arc=0pt,
    left=8pt, right=8pt, top=8pt, bottom=6pt,
    borderline north={2pt}{0pt}{fachfarbe},
    before skip=8pt, after skip=8pt,
    #1
}

\newtcolorbox{formelkasten}{
    colback=white,
    colframe=fachfarbe,
    boxrule=1.5pt,
    arc=0pt,
    left=8pt, right=8pt, top=6pt, bottom=6pt,
    before skip=6pt, after skip=6pt
}

\newtcolorbox{hinweiskasten}{
    colback=white,
    colframe=white,
    boxrule=0pt,
    arc=0pt,
    left=8pt, right=4pt, top=4pt, bottom=4pt,
    borderline west={3pt}{0pt}{fachfarbe}
}

% Befehle
\newcommand{\luecke}[1]{\underline{\hspace{#1}}}
\newcommand{\punktebox}[1]{\hfill\small\textcolor{mittelgrau}{[\hspace{0.8em}/#1]}}
\newcommand{\stern}[1]{\ifzeigesternchen\textsuperscript{\textcolor{fachfarbe}{(#1)}}\fi}

\newcommand{\aufgabe}[2]{%
    \needspace{4\baselineskip}%
    \vspace{8pt}%
    \noindent\textbf{\textcolor{fachfarbe}{Aufgabe #1}} #2%
    \vspace{4pt}%
    \nopagebreak[4]%
}

\newcommand{\operator}[1]{\textbf{#1}}

\newlist{teil}{enumerate}{2}
\setlist[teil,1]{label=\alph*), leftmargin=1.8em, itemsep=4pt, parsep=0pt, topsep=4pt}

\newcommand{\antwortzeilen}[1]{%
    \par\vspace{4pt}%
    \foreach \n in {1,...,#1}{%
        \noindent\rule{\linewidth}{0.3pt}\vspace{8pt}%
    }%
}

\setlength{\parindent}{0pt}
\setlength{\parskip}{3pt}

\begin{document}

% ═══════════════════════════════════════════════════════════════
% HEADER
% ═══════════════════════════════════════════════════════════════

\begin{center}
    {\Large\bfseries\textcolor{fachfarbe}{\thema}}\\[0.2em]
\end{center}
\vspace{-0.5em}
\hrule
\vspace{0.8em}

% ═══════════════════════════════════════════════════════════════
% KASTEN 1: WASSERMODELL
% ═══════════════════════════════════════════════════════════════

\begin{merkkasten}
\textbf{Kasten 1: Die Wasser-Analogie}

\vspace{0.3em}
Elektrischer Strom verhält sich ähnlich wie fließendes Wasser:

\vspace{0.5em}
\begin{center}
\begin{tabular}{|>{\raggedright\arraybackslash}p{4.5cm}|p{5cm}|}
\hline
\textbf{Wasserkreislauf} & \textbf{Stromkreis} \\
\hline
Wasserdruck (durch Pumpe) & \luecke{3.5cm} \\[0.6em]
\hline
Wassermenge pro Sekunde & \luecke{3.5cm} \\[0.6em]
\hline
Engstelle im Rohr & \luecke{3.5cm} \\[0.6em]
\hline
\end{tabular}
\end{center}
\end{merkkasten}

% ═══════════════════════════════════════════════════════════════
% AUFGABE 1: MESSWERTE (NUR U UND I!)
% ═══════════════════════════════════════════════════════════════

\aufgabe{1}{\stern{*} – Messwerte aufnehmen} \punktebox{6}

\operator{Nimm} mithilfe der Simulation Messwerte auf. \operator{Trage} Spannung und Stromstärke in die Tabelle ein.

\vspace{0.5em}
\begin{center}
\begin{tabular}{|c|c||c|}
\hline
\textbf{Spannung $U$ in V} & \textbf{Stromstärke $I$ in A} & \textbf{Widerstand $R$ in $\Omega$} \\
& & \textit{\footnotesize (später ergänzen!)} \\
\hline
2 & & \\[0.7em]
\hline
4 & & \\[0.7em]
\hline
6 & & \\[0.7em]
\hline
8 & & \\[0.7em]
\hline
10 & & \\[0.7em]
\hline
12 & & \\[0.7em]
\hline
\end{tabular}
\end{center}

% ═══════════════════════════════════════════════════════════════
% AUFGABE 2: KENNLINIE ZEICHNEN
% ═══════════════════════════════════════════════════════════════

\aufgabe{2}{\stern{*} – Kennlinie zeichnen} \punktebox{8}

\begin{teil}
    \item \operator{Trage} deine Messpunkte in das Koordinatensystem ein.
    \item \operator{Zeichne} eine \textbf{Ausgleichskurve} durch die Punkte (mit Lineal!).
    \item \operator{Verlängere} die Kurve bis zum Rand des Koordinatensystems.
\end{teil}

\vspace{0.3em}
\begin{center}
\begin{tikzpicture}[scale=0.8]
    % Achsen
    \draw[thick, ->] (0,0) -- (14,0) node[right] {$U$ in V};
    \draw[thick, ->] (0,0) -- (0,8.5) node[above] {$I$ in A};

    % Gitternetz
    \draw[gray!30, very thin] (0,0) grid[step=1] (13,7);

    % Beschriftung x-Achse (0 bis 12 V)
    \foreach \x in {0,2,4,6,8,10,12}
        \draw (\x,0) -- (\x,-0.15) node[below] {\small\x};

    % Beschriftung y-Achse (0 bis 1,4 A, Schritt 0,2)
    % 1 Grid-Einheit = 0,2 A, also y=5 entspricht 1,0 A
    \foreach \y/\ylabel in {0/0, 1/0{,}2, 2/0{,}4, 3/0{,}6, 4/0{,}8, 5/1{,}0, 6/1{,}2, 7/1{,}4}
        \draw (0,\y) -- (-0.15,\y) node[left] {\small\ylabel};
\end{tikzpicture}
\end{center}

\vspace{0.3em}
\textit{Was fällt dir auf?} Die Messpunkte liegen auf einer \luecke{3cm} durch den \luecke{2.5cm}.

% ═══════════════════════════════════════════════════════════════
% SEITE 2
% ═══════════════════════════════════════════════════════════════

\newpage

% ═══════════════════════════════════════════════════════════════
% KASTEN 2: OHMSCHES GESETZ (JETZT ERST!)
% ═══════════════════════════════════════════════════════════════

\begin{merkkasten}
\textbf{Kasten 2: Das Ohmsche Gesetz}

\vspace{0.4em}
\textbf{Deine Entdeckung:} Die Kennlinie ist eine Gerade durch den Ursprung!

\vspace{0.3em}
Das bedeutet: Die \luecke{3cm} ist \textbf{proportional} zur \luecke{3cm}.

\vspace{0.3em}
In Formelzeichen: \hspace{1em} \fbox{\large $I \sim U$} \hspace{1em} (bei konstantem Widerstand)

\vspace{0.6em}
\hrule
\vspace{0.6em}

\textbf{Das Ohmsche Gesetz} gibt dieser Proportionalitätskonstante einen Namen:

\begin{formelkasten}
\begin{center}
\large
$R = \dfrac{U}{I}$ \hspace{2cm} \textbf{Widerstand} $=$ \textbf{Spannung} geteilt durch \textbf{Stromstärke}
\end{center}
\end{formelkasten}

\vspace{0.3em}
\textbf{Einheit:} $[R] = 1\,\Omega$ (Ohm) $= 1\,\dfrac{\text{V}}{\text{A}}$ (Volt pro Ampere)

\vspace{0.6em}
\textbf{Die drei Formen} (durch Umstellen):

\vspace{0.3em}
\begin{center}
\begin{tabular}{|c|c|c|}
\hline
\textbf{Widerstand} & \textbf{Spannung} & \textbf{Stromstärke} \\[0.3em]
\hline
& & \\[0.3em]
$R = \dfrac{U}{I}$ & $U = $ \luecke{2cm} & $I = $ \luecke{2cm} \\[0.8em]
\hline
\end{tabular}
\end{center}

\vspace{0.5em}
\begin{minipage}[t]{0.5\linewidth}
\textbf{Formeldreieck} (Merkhilfe):

\begin{center}
\begin{tikzpicture}[scale=0.6]
    \draw[thick, fachfarbe] (0,0) -- (4,0) -- (2,3) -- cycle;
    \draw[fachfarbe] (0.5,1.5) -- (3.5,1.5);
    \node at (2,2.2) {\large $U$};
    \node at (0.9,0.6) {\large $R$};
    \node at (3.1,0.6) {\large $I$};
    \node at (2,0.6) {\large $\cdot$};
\end{tikzpicture}
\end{center}
\end{minipage}%
\begin{minipage}[t]{0.5\linewidth}
\small
\textit{Anwendung:}\\[0.3em]
Was du suchst, deckst du ab!\\[0.2em]
• $U$ abdecken $\rightarrow$ $U = R \cdot I$\\[0.1em]
• $R$ abdecken $\rightarrow$ $R = U / I$\\[0.1em]
• $I$ abdecken $\rightarrow$ $I = U / R$
\end{minipage}
\end{merkkasten}

% ═══════════════════════════════════════════════════════════════
% AUFGABE 3: R BERECHNEN (JETZT!)
% ═══════════════════════════════════════════════════════════════

\aufgabe{3}{\stern{*} – Widerstand berechnen} \punktebox{4}

\operator{Berechne} jetzt den Widerstand $R$ für alle deine Messwerte aus Aufgabe 1.

\vspace{0.3em}
\begin{hinweiskasten}
\textit{Gehe zurück zu Aufgabe 1 und fülle die rechte Spalte aus!}
\end{hinweiskasten}

\vspace{0.3em}
Was fällt dir auf? \luecke{10cm}

% ═══════════════════════════════════════════════════════════════
% AUFGABE 4: BERECHNUNGEN ÜBEN
% ═══════════════════════════════════════════════════════════════

\aufgabe{4}{\stern{*}/\stern{**} – Berechnungen mit dem Ohmschen Gesetz} \punktebox{12}

\begin{teil}
    \item \stern{*} Gegeben: $U = 12\,\text{V}$, $I = 3\,\text{A}$. \operator{Berechne} $R$. \hfill (2 BE)

    \vspace{0.3em}
    Formel: \luecke{2cm} \hspace{0.5cm} Rechnung: \luecke{4cm} \hspace{0.5cm} $R = $ \luecke{2cm}

    \vspace{0.6em}
    \item \stern{*} Gegeben: $R = 5\,\Omega$, $I = 2\,\text{A}$. \operator{Berechne} $U$. \hfill (2 BE)

    \vspace{0.3em}
    Formel: \luecke{2cm} \hspace{0.5cm} Rechnung: \luecke{4cm} \hspace{0.5cm} $U = $ \luecke{2cm}

    \vspace{0.6em}
    \item \stern{*} Gegeben: $U = 24\,\text{V}$, $R = 8\,\Omega$. \operator{Berechne} $I$. \hfill (2 BE)

    \vspace{0.3em}
    Formel: \luecke{2cm} \hspace{0.5cm} Rechnung: \luecke{4cm} \hspace{0.5cm} $I = $ \luecke{2cm}

    \vspace{0.6em}
    \item \stern{**} Ein Widerstand hat $R = 50\,\Omega$. Bei welcher Spannung fließen $I = 0{,}4\,\text{A}$? \hfill (3 BE)

    \antwortzeilen{2}

    \item \stern{**} Gegeben: $R = 100\,\Omega$, $U = 20\,\text{V}$. Berechne $I$ \textbf{in mA}. \hfill (3 BE)

    \antwortzeilen{2}
\end{teil}

% ═══════════════════════════════════════════════════════════════
% SEITE 3
% ═══════════════════════════════════════════════════════════════

\newpage

% ═══════════════════════════════════════════════════════════════
% AUFGABE 5: GLÜHLAMPE
% ═══════════════════════════════════════════════════════════════

\aufgabe{5}{\stern{**} – Die Glühlampe: Ein besonderer Widerstand} \punktebox{11}

\begin{teil}
    \item \operator{Nimm} mit der zweiten Simulation Messwerte für eine \textbf{Glühlampe} auf. \operator{Berechne} $R = U/I$. \hfill (4 BE)

    \vspace{0.3em}
    \begin{center}
    \small
    \begin{tabular}{|c|c|c|c|c|c|c|}
    \hline
    $U$ in V & 2 & 4 & 6 & 8 & 10 & 12 \\
    \hline
    $I$ in A & & & & & & \\[0.5em]
    \hline
    $R$ in $\Omega$ & & & & & & \\[0.5em]
    \hline
    \end{tabular}
    \end{center}

    \vspace{0.3em}
    Was fällt dir bei den $R$-Werten auf? \luecke{6cm}

    \vspace{0.5em}
    \item \operator{Zeichne} die Kennlinie der Glühlampe in \textbf{einer anderen Farbe} in das Diagramm auf Seite 1 ein. \hfill (2 BE)

    \vspace{0.5em}
    \item \operator{Vergleiche} die beiden Kennlinien. Was ist anders? \hfill (2 BE)

    \antwortzeilen{2}

    \vspace{0.3em}
    \item \operator{Erkläre}, warum die Kennlinie der Glühlampe gekrümmt ist. \hfill (3 BE)

    \vspace{0.3em}
    \textit{Fülle die Lücken:}

    \vspace{0.4em}
    Bei höherer Spannung fließt mehr \luecke{2.5cm}. Die Elektronen stoßen häufiger mit

    \vspace{0.3em}
    den \luecke{3cm}. Dadurch beginnen diese stärker zu \luecke{2.5cm}.

    \vspace{0.3em}
    Der Draht wird \luecke{2cm}. Je heißer der Draht, desto \luecke{2cm} ist der

    \vspace{0.3em}
    Widerstand. Deshalb steigt der Strom \luecke{2.5cm} stark als erwartet.

    \vspace{0.4em}
    \hfill\textit{\small (Wörter: Strom, Atomrümpfen, schwingen, heiß, größer, weniger)}
\end{teil}

% ═══════════════════════════════════════════════════════════════
% AUFGABE 6: TRANSFER
% ═══════════════════════════════════════════════════════════════

\aufgabe{6}{\stern{***} – Anwendungen} \punktebox{6}

\begin{teil}
    \item Ein Föhn ist mit $230\,\text{V}$ und $10\,\text{A}$ beschriftet. \operator{Berechne} $R$. \hfill (2 BE)

    \antwortzeilen{2}

    \item Verlängerungskabel werden manchmal warm. \operator{Erkläre} warum. \hfill (2 BE)

    \antwortzeilen{2}

    \item Lisa berechnet $R = 8\,\Omega$. Richtig wäre $R = 4\,\Omega$. \operator{Erkläre} ihren Fehler. \hfill (2 BE)

    \antwortzeilen{2}
\end{teil}

% ═══════════════════════════════════════════════════════════════
% KASTEN 3: ZUSAMMENFASSUNG
% ═══════════════════════════════════════════════════════════════

\vspace{0.3em}
\begin{merkkasten}
\textbf{Kasten 3: Zusammenfassung}

\vspace{0.3em}
\begin{itemize}[leftmargin=1.5em, itemsep=2pt]
    \item Ohmsches Gesetz: \luecke{2cm} $\sim$ \luecke{2cm} (Proportionalität)
    \item Formeln: $R =$ \luecke{1.5cm}, \hspace{0.2cm} $U =$ \luecke{1.5cm}, \hspace{0.2cm} $I =$ \luecke{1.5cm}
    \item Einheit: $[R] = 1\,\Omega = $ \luecke{2cm}
    \item Ohmsche Widerstände: \luecke{2.5cm} Kennlinie (R konstant)
    \item Glühlampen: \luecke{2.5cm} Kennlinie (R steigt mit \luecke{2cm})
\end{itemize}
\end{merkkasten}

\end{document}
